\documentclass[10pt,a4paper]{amsart}
\usepackage[margin=19mm]{geometry}
\usepackage{amsmath,amssymb,mathtools}
\usepackage[scr=rsfs,cal=euler]{mathalfa}
\usepackage{xfrac}
\usepackage[unicode,hidelinks,bookmarks=false]{hyperref}
\newcommand{\ie}{i.e.\ }
\newcommand{\cf}{cf.\ }
\newcommand{\resp}{respectively\ }
\newcommand{\Subsection}{\subsection}
\providecommand{\nobreakdash}{\nobreak\hbox{-}}
\setlength{\emergencystretch}{2em}
\multlinegap0pt
\usepackage{mdframed}
\usepackage{mdframed}
\usepackage{mdframed}
\usepackage{mdframed}
\usepackage{amssymb}
\usepackage{mathtools}
\usepackage{bm}
\usepackage[normalem]{ulem}
\usepackage{enumerate}
\usepackage{xcolor}
\usepackage{tikz}
\usepackage{mdframed}
\usepackage{cancel}
\newtheorem{lemma}{Lemma}
\newcommand{\E}{\mbox{${\mathbb E}$}}
\newcommand{\FF}{\mbox{${\mathcal F}$}}
\newcommand{\I}{\mbox{${\mathbb I}$}}
\title{Tampered Memory Elephant Random Walk on one-dimensional integer lattice}
\author{Vinita Mulay, Neeraja Sahasrabudhe, Debleena Thacker}
\date{Source: arXiv:2607.28614v1 (CC BY 4.0)}
\begin{document}
\maketitle

\noindent\emph{Source and licence.} Tampered Memory Elephant Random Walk on one-dimensional integer lattice. Version arXiv:2607.28614v1 (CC BY 4.0), distributed by arXiv under the Creative Commons Attribution 4.0 International licence, http://creativecommons.org/licenses/by/4.0/, as recorded in the arXiv licence metadata for this exact version and shown on the abstract page https://arxiv.org/abs/2607.28614v1. Full licence notice: \texttt{licenses/arxiv-2607.28614-CC-BY-4.0.txt}.\par\smallskip

\noindent\textit{Editorial scope.} Counted unit: the article's lemma lem:crossterms (source0.tex lines 1287--1291). The article's complete original proof of it is delivered with it (source0.tex lines 1292--1405, one \texttt{proof} environment of 114 lines). No mathematics is written, repaired or re-derived: the statement and the proof are the article's own lines, in the article's own notation, and no retained line is substituted or re-flowed.  No further source-local context is needed: every cross-reference of every retained line resolves inside the two delivered spans.  Nothing was omitted from any retained span, so the package carries no omission record.  No retained line contains a citation, so the wrapper carries no bibliography and no bibliography entry.  No retained line contains a figure environment, an image-inclusion command or drawing code, so no image is delivered and the package declares no asset dependency.  Editorial formatting only, in addition: an AMS \texttt{amsart} preamble that restates the article's own declarations and macros used by the retained lines, namely the environments \texttt{lemma} on separate counters, together with the standard packages those lines need.  The retained environments print in this document as Lemma 1 in order of appearance, whereas the article prints the same items with the numbering of its own full document.  The complete native source of the article as distributed in the arXiv e-print for this exact version is bundled unchanged as \texttt{sources/206429/source0.tex}.
\begin{lemma} \label{lem:crossterms}
For all $i \neq j \in \{0, 1, \dots, \mu-1 \}$
\[ \lim_{n \to \infty} \E\left[ X_{\mu n+i} X_{\mu n+j} \vert \FF_{\mu n} \right] = \eta^2, \]
where $\eta = \frac{1}{\mu}((2p-1){k}s^* +(2\lambda-1)(\mu-k))$. 
\end{lemma}

\begin{proof}
    Consider $i, j$ such that $1 \leq i < j \leq \mu$. We have,
    \[ \E[X_{\mu n+i}X_{\mu n+j}\vert \FF_{\mu n}] = \E \left[ X_{\mu n+i} \E[X_{\mu n+j}\vert  \FF_{\mu n + j-1} ] \Bigm| \FF_{\mu n} \right]. \]
    Since, \[\E[X_{\mu n+j}\vert \FF_{\mu n + j-1}] = (2p-1) \frac{S_{\mu n+j-1}^{D^c}}{\mu n+j-1} + (2\lambda-1) \frac{|D_{\mu n+j-1}|}{\mu n+j-1}.\]
    and, $S_{\mu n+j-1}^{D^c} = S_{\mu n}^{D^c} + \sum_{r=1}^{j-1} \I_{\{\mu n+r\in D^c\}} X_{\mu n+r}$. 
It follows that
\begin{align*}
E[X_{\mu n+j}\mid  \FF_{\mu n + j-1} ]
&=
(2p-1)
\frac{
S_{\mu n}^{D^c}
+
\sum_{r=1}^{j-1}
\I_{\{\mu n+r\in D^c\}}
X_{\mu n+r}
}
{\mu n+j-1}
\\
&\qquad
+
(2\lambda-1)
\frac{|D_{\mu n+j-1}|}{\mu n+j-1}.
\end{align*}
Substituting this into the tower property gives
\begin{align*}
\E[X_{\mu n+i}X_{\mu n+j}\mid \FF_{\mu n}]
&=
\frac{2p-1}{\mu n+j-1}
E\left[
X_{\mu n+i}
\left(
S_{\mu n}^{D^c}
+
\sum_{r=1}^{j-1}
\I_{\{\mu n+r\in D^c\}}
X_{\mu n+r}
\right)
\Bigm|
\FF_{\mu n}
\right]
\\
&\qquad
+
(2\lambda-1)
\frac{|D_{\mu n+j-1}|}{\mu n+j-1}
\E[X_{\mu n+i}\mid \FF_{\mu n}].
\end{align*}
Expanding the first expectation,
\begin{align*}
\E[X_{\mu n+i}X_{\mu n+j}\mid \FF_{\mu n}]
&= \frac{2p-1}{\mu n+j-1}
\sum_{r=1}^{j-1}
\I_{\{\mu n+r\in D^c\}}
\E[X_{\mu n+i}X_{\mu n+r}\mid \FF_{\mu n}]
\\
& + \frac{2p-1}{\mu n+j-1}
S_{\mu n}^{D^c}
E[X_{\mu n+i}\mid \FF_{\mu n}] +
(2\lambda-1)
\frac{|D_{\mu n+j-1}|}{\mu n+j-1}
\E[X_{\mu n+i}\mid \FF_{\mu n}].
\end{align*}

Since $j\le \mu$ is fixed, the sum in the first term contains at most $\mu$ terms. Since $\vert \E[X_{\mu n+i}X_{\mu n+r}\mid \FF_{\mu n}] \vert
\leq 1$, for some constant $C>0$
\begin{equation*} 
\left|
\frac{2p-1}{\mu n+j-1}
\sum_{r=1}^{j-1}
\I_{\{\mu n+r\in D^c\}}
\E[X_{\mu n+i}X_{\mu n+r}\mid \FF_{\mu n}]
\right|
\le
\frac{C}{n}    
\end{equation*}
Thus,
\begin{align*}
\E[X_{\mu n+i}X_{\mu n+j}\mid \FF_{\mu n}]
&=
\Bigg(
(2p-1)\frac{S_{\mu n}^{D^c}}{\mu n+j-1}
+ (2\lambda-1)\frac{|D_{\mu n+j-1}|}{\mu n+j-1}
\Bigg)
\E[X_{\mu n+i}\mid \FF_{\mu n}] +o(1).
\end{align*}
Since $i$ and $j$ are fixed,
\[
\frac{S_{\mu n}^{D^c}}{\mu n+j-1} = \frac{S_{\mu n}^{D^c}}{\mu n} +o(1)
\quad \text{ and } \quad 
\frac{|D_{\mu n+j-1}|}{\mu n+j-1} =
\frac{\mu-k}{\mu}+o(1).
\]
Similarly,
\[ 
\E[X_{\mu n+i}\mid \FF_{\mu n}] = (2p-1)\frac{S_{\mu n}^{D^c}}{\mu n} + (2\lambda-1)\frac{\mu-k}{\mu}
+o(1).
\]
Using $\frac{S_{\mu n}^{D^c}}{n} \longrightarrow ks^*$, 
we obtain $\E[X_{\mu n+i}\mid \FF_{\mu n}]
\longrightarrow \eta$,
where $\eta = \frac{(2p-1)ks^*+(2\lambda-1)(\mu-k)}{\mu}$. Similarly, 
\[
(2p-1)\frac{S_{\mu n}^{D^c}}{\mu n+j-1}
+ (2\lambda-1)\frac{|D_{\mu n+j-1}|}{\mu n+j-1} \longrightarrow \eta.
\]
Therefore,
\[
\E[X_{\mu n+i}X_{\mu n+j}\mid \FF_{\mu n}]
\longrightarrow
\eta^2,
\]
which completes the proof.
\end{proof}

\end{document}
